El costo del proceso actual de gestión de materiales no se limita al tiempo que invierte el equipo de gestores: se traslada a toda la cadena de procesos que depende de que un material quede correctamente registrado en el sistema. Mientras una solicitud permanece en gestión (en promedio 3.3 horas dentro del rol del gestor, sobre un volumen de 4.0 solicitudes diarias), los procesos que suceden, como la emisión de una orden de compra o la ejecución de una tarea de mantenimiento, quedan detenidos. Reducir ese tiempo no mejora únicamente la productividad del gestor: mejora la velocidad de respuesta de la organización.

El impacto se agrava cuando el resultado de la gestión no es correcto. Un catálogo que ya arrastra 2,089 duplicados exactos, otros 3,649 registros sin la estructura que permitiría detectarlos por búsqueda exacta y materiales asignados a categorías que no corresponden a su naturaleza, no solo es un problema de orden interno: distorsiona los reportes de consumo, complica la planificación de compras por tipo de material y reduce la trazabilidad de los activos para cualquier área que consulte esa información después. Cada día que el proceso continúa sin apoyo analítico, el catálogo acumula más inconsistencias, y el costo de corregirlas retroactivamente crece.

Atender este problema requiere visibilidad sobre cómo se está comportando el proceso hoy, y una forma de intervenir directamente en el momento en que ocurre el error. Un tablero que solo reporte el problema, sin cambiar la forma en que se crea el material, no resuelve la causa; y automatizar la creación sin poder medir su efecto no permite sostener la mejora en el tiempo. De ahí que la solución combine ambas capas: un modelo que asiste al gestor en el instante de la decisión, y un mecanismo de medición que permite verificar, con datos, si esa asistencia efectivamente está reduciendo el tiempo de gestión y mejorando la calidad del catálogo.

La información necesaria para construir y validar esta solución es de acceso directo dentro de la unidad de negocio, y la infraestructura tecnológica requerida se apoya íntegramente en herramientas de código abierto, lo que hace sostenible su operación más allá del período de desarrollo, sin inversión adicional significativa en licenciamiento.

Finalmente, el proyecto se justifica también en términos financieros: la reducción del tiempo de gestión y la prevención de duplicados se traduce en un ahorro cuantificable de horas-gestor, cifra que se desarrolla en detalle en la sección de Resultados de este documento.